\scanpage{13}
\deleted{Considérons un point $\xi\in S(k)$, \inserted{$k$ corps de car. $0$,} il définit un foncteur fibre $H^{\unsure{**|*}}_\xi$ (Hodge) et $H^*_\xi$ (De Rham) : \unsure{à valeurs} dans la \uline{catégorie} des vectoriels sur $k$.}

Cas d'un \uline{tore} $T$ sur $\mathbb Q$.

Défini pr un $\pi$-module admissible, libre d.t.f. sur $\mathbb Z$, soit $M$, ($\pi=\pi_1(\mathbb Q,\mathbb C)$)

On aura $\varepsilon\in M^\pi$, $\unsure{i|j}\in\check M^\pi$, $\langle i,\varepsilon\rangle=2$.

Les modules sur $T$ correspondent aux éléments invariants pr $\pi$ de $\mathbb Z(M)^+$ (Comb. lin. à coeff. $\geq0$)

\deleted{Donnons nous \unsure{un}}

La donnée d'une structure \deleted{\unsure{positive}} \inserted{de polarisation} ne dépend que du tore réel $T_{\mathbb R}$. Ce dernier est de la forme \deleted{$T_{\mathbb R}\simeq U\times S$,} \inserted{ext. d'un $S$ splitté pr} $U$ compact, \inserted{i.e. $U(\mathbb R)$ cpact,} \deleted{$S$ \unsure{splitté}} et $\varepsilon$ est néc. trivial sur $U$. Donc on en déduit des structures $\varepsilon_S$ et $j_S$ sur $S$, et évidemment $S$ \uline{hérite} de structure de polarisation. J'affirme qu'on doit avoir $\dim S=1$, en d'autres termes, que \uline{$\operatorname{Ker}\varepsilon$ est compact}.

\deleted{En effet, si on avait $\dim S\geq2$, on pourrait \unsure{peut} \deleted{trouver un} \unsure{un caractère} \unsure{Choisissons} \unsure{par $\pi$} \unsure{les} \unsure{sur droite} \unsure{caractère} de $\operatorname{Ker}\varepsilon$) que la \unsure{dim} \unsure{est deux}. \deleted{Soit} On \unsure{peut} écrire
\[
S=\mathbb G_{m\mathbb R}\times\mathbb G_{m\mathbb R},\quad\text{avec }i(\lambda)=(\unsure{0|1},\lambda^\alpha),\quad\varepsilon(\lambda,\mu)=\mu^\beta
\]
où $\alpha,\beta$ entiers $\geq1$, $\alpha\beta=2$ (donc $\alpha=1$, $\beta=2$, ou $\alpha=2$, $\beta=1$).}

\mspage{12}
